\documentclass[twocolumn]{article}
\usepackage{amsmath}
\usepackage{amssymb}
\usepackage{algorithm}
\usepackage{algpseudocode}
\usepackage{booktabs}
\usepackage{parskip}
\usepackage{float}
\usepackage{geometry}
\usepackage{pifont}
\usepackage{xcolor}
\usepackage[font={small,it}, labelfont={bf,it}, labelsep=period]{caption}
\geometry{margin=0.75in, columnsep=0.25in}
\newcommand{\cmark}{\ding{51}}
\newcommand{\xmark}{\ding{55}}

% Title and author defined before \begin{document} for twocolumn
\title{\textbf{Before the Search Begins: Generating Function-Informed Initialization for Tight Constrained Binary Optimization}}
\author{Muhammad Marzouk Baig\\
{\small MS Candidate, Khoury College of Computer Sciences}\\
{\small Northeastern University, Portland, Maine}\\
{\small \texttt{baig.muham@northeastern.edu}}}
\date{}

\begin{document}

\twocolumn[
  \maketitle
  \begin{center}
  \begin{minipage}{0.85\textwidth}
  \small
  \begin{abstract}
We present a generating function-informed initialization framework for genetic 
algorithms (GAs) applied to the Multidimensional Knapsack Problem (MKP). 
Standard initialization---setting each bit to 1 with probability 0.5 
independently---ignores the feasibility structure of the search space entirely. 
On tight instances this produces initial populations with zero feasible 
solutions, preventing any search from proceeding under constraint-respecting 
fitness models. The framework is grounded in the combinatorial generating 
function tradition, applied here to constraint feasibility rather than 
partition enumeration: for each item $j$, we derive a sampling probability 
$p_j$ from sub-generating functions that measure how much feasible space 
is lost when item $j$ is included versus excluded. The framework is purely 
constraint-driven: it requires no objective function, no LP relaxation, and 
no prior search. A GPU-accelerated Monte Carlo approximation provides 
estimates accurate to within $0.03\%$ in 5--8 seconds. On tight instances 
($\alpha = 0.25$), where uniform initialization produces zero feasible 
solutions, our method achieves $52.5\%$ initial feasibility and converges 
to $96.5\%$ of the reference best under a death penalty model---a regime 
where the uniformly initialized GA fails entirely. Adding a lightweight 
iterated local search post-processing step raises mean quality to $99.3\%$. 
The framework is designed for tight-to-medium instances ($\alpha \leq 0.50$) 
where feasibility is genuinely scarce; on loose instances it degenerates 
automatically to uniform initialization.
  \end{abstract}
  \end{minipage}
  \end{center}
  \vspace{3em}
]

% ============================================================
\section{Introduction}
% ============================================================

\subsection{Motivation}

Population initialization is the first decision a genetic algorithm makes, 
yet it is almost universally treated as a formality. The standard approach---sampling 
each bit independently with probability $0.5$---is convenient and assumption-free, 
but it is provably oblivious to the structure of the problem being solved. 
For unconstrained or loosely constrained problems this is harmless. For 
tightly constrained combinatorial problems, it can be catastrophic.

\textbf{The generating function connection.}
Generating functions are a classical tool in combinatorics for encoding 
subset enumeration in compact algebraic form. The idea goes back at least 
to Euler and reached remarkable depth in the work of Hardy and 
Ramanujan~\cite{hardyramanujan1918} on integer partitions: the generating 
function $\prod_{k=1}^{\infty}(1-x^k)^{-1}$ encodes, in the coefficient of 
$x^n$, the number of ways to write $n$ as a sum of positive integers. 
The same algebraic principle applies to our setting. For items with integer 
weights $w_1, \ldots, w_n$, the product $\prod_j (1 + x^{w_j})$ encodes, 
in the coefficient of $x^k$, the number of subsets with total weight 
exactly $k$. This maps naturally onto the feasibility question for 
constrained binary optimization: by factoring out one item at a time and 
evaluating the remainder against each capacity threshold, we can 
ask---before any search begins---how much of the feasible space survives 
when item $j$ is included versus excluded. To our knowledge, this 
connection has not previously been exploited to bias GA initialization. 
The framework we develop shows that it yields a principled, interpretable, 
and provably self-calibrating initialization method.

\textbf{The tight instance failure mode.}
Consider the MKP with a tightness ratio $\alpha = W_i / \sum_j w_{ji}$ of 
$0.25$: each capacity is set to $25\%$ of the total item weight across 
that constraint. On such instances the fraction of all binary strings that 
are feasible is tiny. Uniform random initialization samples from the full 
$\{0,1\}^n$ space without any awareness of where feasible solutions lie; 
in practice, every solution in an initial population of 50 is infeasible. 
Under a death penalty fitness model---which assigns zero fitness to 
infeasible solutions---the entire initial population is non-competing and 
selection cannot operate. The GA is dead on arrival. This failure mode is 
not an artifact of a poorly chosen penalty function or an implementation 
bug. It is a structural consequence of initializing without any knowledge 
of where feasible solutions live.

\textbf{The scope of repair-based approaches.}
The most effective existing technique for handling infeasibility in MKP GAs 
is the deterministic repair operator of Chu and 
Beasley~\cite{chu1998genetic}, which uses an efficiency ranking 
$u_j = v_j / \sum_i (w_{ji} / W_i)$ to greedily add or remove items until 
feasibility is restored. This operator is computationally efficient and 
empirically effective on standard benchmarks. However, it depends entirely 
on objective values: the ranking cannot be computed when objective values 
are absent (pure constraint satisfaction), unreliable (noisy or estimated 
objectives), or misleading (adversarial instance structure). In these 
settings, constraint-only initialization methods become the primary 
available tool for guiding early search.

\subsection{Contributions}

The specific contributions of this paper are:

\begin{enumerate}
    \item A mathematical framework that uses sub-generating functions over 
    order-1 schemas to derive per-item feasibility densities 
    $\rho_j^1$ and $\rho_j^0$ for any binary constrained optimization 
    instance (Section~\ref{sec:framework}).

    \item A sampling rule $p_j = \rho_j^1 / (\rho_j^1 + \rho_j^0)$ 
    that recovers uniform initialization when no constraint structure 
    is present ($p_j = 0.5$) and biases away from infeasibility-inducing 
    items in proportion to the feasibility cost they impose 
    (Section~\ref{sec:framework}).

    \item A GPU-accelerated Monte Carlo approximation algorithm that 
    estimates $\rho_j^1$ and $\rho_j^0$ to within $0.03\%$ in 5--8 
    seconds for $n = 100$, $m = 3$, running in $O(n \cdot M \cdot m)$ 
    time with $M = 10^7$ samples (Section~\ref{sec:algorithm}).

    \item Empirical results on MKP instances at three tightness levels 
    ($\alpha \in \{0.25, 0.50, 0.75\}$) showing that GF-biased 
    initialization enables search under a death penalty model in regimes 
    where uniform initialization completely fails, while matching uniform 
    performance on easier instances. We also show that the framework 
    composes naturally with standard post-processing techniques 
    (Section~\ref{sec:results}).
\end{enumerate}

\subsection{Paper Organization}

Section~\ref{sec:background} reviews the MKP, the Chu--Beasley GA, 
schema theory, and generating functions. Section~\ref{sec:framework} 
develops the mathematical framework. Section~\ref{sec:algorithm} presents 
the exact and Monte Carlo algorithms along with a complexity analysis. 
Section~\ref{sec:examples} provides worked examples on small instances. 
Section~\ref{sec:experiments} describes the experimental setup. 
Section~\ref{sec:results} presents results. Section~\ref{sec:discussion} 
discusses limitations and the broader applicability of the framework. 
Section~\ref{sec:conclusion} concludes.

% ============================================================
\section{Background}
\label{sec:background}
% ============================================================

\subsection{The Multidimensional Knapsack Problem}

The MKP generalizes the classical 0-1 knapsack problem to $m$ constraints. 
Given $n$ items each with value $v_j$ and weight $w_{ji}$ for constraint 
$i$, and $m$ capacity constraints $W_i$, the MKP is:

\begin{flushleft}
\hspace{2em}$\text{maximize} \quad \displaystyle\sum_{j=1}^{n} v_j x_j$
\end{flushleft}

\begin{flushleft}
\hspace{2em}$\text{subject to} \quad \displaystyle\sum_{j=1}^{n} w_{ji} x_j 
\leq W_i \quad \forall i = 1, \ldots, m$
\end{flushleft}

\begin{flushleft}
\hspace{2em}$x_j \in \{0, 1\} \quad \forall j = 1, \ldots, n$
\end{flushleft}

The MKP is NP-hard and is a canonical benchmark for metaheuristics. 
Instances are commonly parameterized by the \textit{tightness ratio} 
$\alpha = W_i / \sum_j w_{ji}$, which controls the fraction of the total 
item weight that fits in each constraint. Standard benchmark instances 
are available from the OR-Library~\cite{beasley1990orlib}.

\subsection{Chu and Beasley Genetic Algorithm}
\label{sec:chundbeasley}

Chu and Beasley~\cite{chu1998genetic} proposed a steady-state GA for the 
MKP that remains the standard baseline. Key components:

\begin{itemize}
    \item \textbf{Initialization:} Each bit $x_j$ set to 1 with probability 
    0.5 independently.
    \item \textbf{Crossover:} Uniform crossover.
    \item \textbf{Mutation:} Bit-flip at rate $1/n$.
    \item \textbf{Selection:} Binary tournament.
    \item \textbf{Repair:} A deterministic operator that iteratively drops 
    items (in ascending order of efficiency $u_j$) until feasibility is 
    restored, then greedily adds items (in descending order of $u_j$) 
    while feasibility is maintained, where:
\end{itemize}

\begin{flushleft}
\hspace{2em}$u_j = \dfrac{v_j}{\displaystyle\sum_{i=1}^{m} w_{ji} / W_i}$
\end{flushleft}

The repair operator is effective on standard benchmarks because item 
efficiency is a reliable proxy for solution quality when values and weights 
are uncorrelated. It is undefined when no objective exists and potentially 
misleading when value-weight correlation is adversarial~\cite{kellerer2004knapsack}. 
Because the repair ranking is fixed, diverse infeasible offspring tend to be 
repaired toward similar feasible solutions, which can contribute to premature 
convergence on harder instances.

\subsection{Schema Theory}

Schema theory, introduced by Holland~\cite{holland1975adaptation}, 
provides a framework for analyzing GA behavior at the level of partial 
solutions. A \textit{schema} is a string of length $n$ over 
$\{0, 1, *\}$, where $*$ is a wildcard matching either value. 
An \textit{order-$k$} schema has $k$ fixed positions. 
The Schema Theorem states that short, low-order, above-average schemata 
are sampled exponentially more frequently across generations. 
Initialization determines the initial representation frequency of all 
schemata; biasing toward high-feasibility order-1 schemata directly 
seeds the population with the building blocks from which feasible 
high-quality solutions will be composed.

\subsection{Generating Functions in Combinatorics}

A generating function is a formal power series whose coefficients encode 
combinatorial quantities. The tradition is ancient: Euler used generating 
functions to analyze integer partitions, and the field reached a landmark 
in the Hardy--Ramanujan asymptotic formula~\cite{hardyramanujan1918} for 
the partition function $p(n)$, derived by analyzing 
$\prod_{k=1}^{\infty}(1-x^k)^{-1}$. For a finite set of items with 
integer weights $w_1, \ldots, w_n$, the analogous product is:

\begin{flushleft}
\hspace{2em}$f(x) = \prod_{j=1}^{n} (1 + x^{w_j})$
\end{flushleft}

The coefficient of $x^k$ in $f(x)$ equals the number of subsets with 
total weight exactly $k$. This follows from the distributive law: each 
factor contributes either $1$ (item excluded) or $x^{w_j}$ (item included), 
so the product enumerates all $2^n$ subsets. Wilf~\cite{wilf1994genfunc} 
provides a comprehensive treatment of generating functions in combinatorics; 
the specific application to constraint feasibility and GA initialization 
developed here is, to our knowledge, new.

% ============================================================
\section{Mathematical Framework}
\label{sec:framework}
% ============================================================

\subsection{Feasibility and the Full Generating Function}

Let there be $n$ items and $m$ constraints. A solution 
$\mathbf{x} \in \{0,1\}^n$ is \textit{feasible} if:

\begin{flushleft}
\hspace{2em}$\displaystyle\sum_{j=1}^{n} w_{ji} x_j \leq W_i 
\quad \forall i = 1, \ldots, m$
\end{flushleft}

For $m$ constraints, the full generating function over all $2^n$ subsets is:

\begin{flushleft}
\hspace{2em}$f(x_1, \ldots, x_m) = \prod_{j=1}^{n} 
\left(1 + x_1^{w_{j1}} x_2^{w_{j2}} \cdots x_m^{w_{jm}}\right)$
\end{flushleft}

A term $x_1^{k_1} \cdots x_m^{k_m}$ in the expanded product corresponds 
to a feasible subset if and only if $k_i \leq W_i$ for all $i$.

\subsection{Sub-Generating Functions and Feasibility Densities}

To analyze feasibility at the schema level, we define the 
\textit{sub-generating function} for item $j$ as the generating function 
over all $n-1$ remaining items:

\begin{flushleft}
\hspace{2em}$g_j(x_1, \ldots, x_m) = \prod_{k \neq j} 
\left(1 + x_1^{w_{k1}} \cdots x_m^{w_{km}}\right)$
\end{flushleft}

This function enumerates all $2^{n-1}$ subsets of items other than $j$. 
The function $g_j$ is the same regardless of whether item $j$ is fixed 
to 0 or 1; the distinction lies solely in the capacity threshold applied 
when counting feasible terms:

\begin{itemize}
    \item When item $j$ is \textbf{included} ($x_j = 1$), its weight 
    consumes capacity, leaving $W_i - w_{ji}$ for the remaining items. 
    The feasibility density is:
\end{itemize}

\begin{flushleft}
\hspace{2em}$\rho_{j}^{\,1} = \dfrac{\#\bigl\{\text{terms in } 
g_j : \text{exponent of } 
x_i \leq W_i - w_{ji} \ \forall i\bigr\}}{2^{n-1}}$
\end{flushleft}

\begin{itemize}
    \item When item $j$ is \textbf{excluded} ($x_j = 0$), full capacity 
    is available. The feasibility density is:
\end{itemize}

\begin{flushleft}
\hspace{2em}$\rho_{j}^{\,0} = \dfrac{\#\bigl\{\text{terms in } 
g_j : \text{exponent of } 
x_i \leq W_i \ \forall i\bigr\}}{2^{n-1}}$
\end{flushleft}

Intuitively, $\rho_j^1$ is the probability that a uniformly random 
solution \textit{containing} item $j$ is feasible; $\rho_j^0$ is the 
same for solutions \textit{not} containing $j$.

\subsection{Order-1 Schemas and the Schema-Level Interpretation}

A schema is a string over $\{0, 1, *\}$ where $*$ is a wildcard. An 
order-1 schema has exactly one fixed position. For $n$ items there are 
$2n$ order-1 schemas---two per item---each containing exactly $2^{n-1}$ 
strings. The sub-generating function $g_j$ corresponds exactly to the 
order-1 schema for item $j$: it enumerates and characterizes the 
feasibility of all $2^{n-1}$ strings that agree with item $j$'s fixed 
value. Restricting the analysis to order-1 structures requires $n$ 
sub-generating function computations (one per item) rather than the $3^n$ 
required for full schema enumeration.

\subsection{Sampling Probability}
\label{sec:samplingprob}

We allocate the initialization probability for bit $j$ proportionally to 
the feasibility signal of each choice. For two mutually exclusive options 
with scores $a$ and $b$, the proportional allocation rule assigns 
probability $a/(a+b)$ to option $a$. Applied to our two feasibility 
densities:

\begin{flushleft}
\hspace{2em}$p_j = \dfrac{\rho_{j}^{\,1}}{\rho_{j}^{\,1} + \rho_{j}^{\,0}}$
\end{flushleft}

This rule satisfies two important properties:
\begin{enumerate}
    \item $p_j + (1 - p_j) = 1$ (the probabilities sum to one by 
    construction).
    \item The ratio of probabilities equals the ratio of feasibility 
    densities: if excluding item $j$ is twice as feasibility-preserving 
    as including it, the algorithm is twice as likely to exclude it.
\end{enumerate}

\noindent \textbf{Boundary cases.}
When $\rho_j^1 = \rho_j^0$, we have $p_j = 0.5$, recovering standard 
uniform initialization. This occurs on loose instances ($\alpha$ near 1), 
where the presence or absence of any single item has negligible effect 
on feasibility. When $\rho_j^1 = 0$, item $j$ is never compatible with 
any completion of the remaining items---$p_j = 0$ and the item is never 
selected. When $\rho_j^0 = 0$, all completions are infeasible regardless 
of whether $j$ is included or not, which cannot occur on any instance 
for which the empty solution is feasible.

% ============================================================
\section{Algorithm}
\label{sec:algorithm}
% ============================================================

\subsection{Exact Computation}
\label{sec:exact}

The exact procedure computes $\rho_j^1$ and $\rho_j^0$ by explicitly 
enumerating all $2^{n-1}$ terms of each sub-generating function. Note 
that the full generating function $f$ itself is never expanded; only the 
$n$ sub-generating functions $g_j$ are required, each of size $2^{n-1}$.

\begin{algorithm}[H]
\caption{GF-Informed Initialization (Exact)}
\label{alg:exact}
\begin{algorithmic}[1]
\Require $n$ items, $m$ constraints, weights $w_{ji}$, 
         capacities $W_i$, population size $P$
\Ensure Initial population of $P$ solutions

\Statex \textbf{Phase 1 --- Precomputation (done once)}
\For{$j = 1$ to $n$}
    \State Expand $g_j(x_1, \ldots, x_m) = 
           \displaystyle\prod_{k \neq j} 
           \left(1 + x_1^{w_{k1}} \cdots x_m^{w_{km}}\right)$
    \State $\rho_{j}^{\,1} \gets$ fraction of terms where 
           exponent of $x_i \leq W_i - w_{ji}$ for all $i$
    \State $\rho_{j}^{\,0} \gets$ fraction of terms where 
           exponent of $x_i \leq W_i$ for all $i$
    \State $p_j \gets \rho_{j}^{\,1} / 
           (\rho_{j}^{\,1} + \rho_{j}^{\,0})$
\EndFor

\Statex \textbf{Phase 2 --- Population Initialization}
\For{$s = 1$ to $P$}
    \For{$j = 1$ to $n$}
        \State $x_j \gets 1$ with probability $p_j$, else $0$
    \EndFor
    \State Add solution $\mathbf{x}$ to population
\EndFor
\State \Return population
\end{algorithmic}
\end{algorithm}

Exact computation requires enumerating $2^{n-1}$ terms per item and is 
feasible only for $n \leq 20$ in practice. Phase 2 is identical to 
standard uniform initialization except each bit $j$ is sampled with 
probability $p_j$ rather than a fixed $0.5$.

\subsection{Monte Carlo Approximation}
\label{sec:montecarlo}

For practical instances ($n = 50$--$500$), exact enumeration is 
intractable. We approximate $\rho_j^1$ and $\rho_j^0$ via Monte Carlo 
sampling. Instead of enumerating all $2^{n-1}$ completions, we draw 
$M$ random binary vectors $\mathbf{s}^{(1)}, \ldots, \mathbf{s}^{(M)}$ 
over the $n-1$ items other than $j$, and count the fraction that are 
feasible under each threshold:

\begin{flushleft}
\hspace{2em}$\hat{\rho}_{j}^{\,1} = \dfrac{1}{M} 
\displaystyle\sum_{k=1}^{M} 
\mathbf{1}\!\left[\displaystyle\sum_{\ell \neq j} w_{\ell i}\, s_\ell^{(k)} 
\leq W_i - w_{ji} \;\; \forall i\right]$
\end{flushleft}

\begin{flushleft}
\hspace{2em}$\hat{\rho}_{j}^{\,0} = \dfrac{1}{M} 
\displaystyle\sum_{k=1}^{M} 
\mathbf{1}\!\left[\displaystyle\sum_{\ell \neq j} w_{\ell i}\, s_\ell^{(k)} 
\leq W_i \;\; \forall i\right]$
\end{flushleft}

\noindent where $s_\ell^{(k)} \overset{\text{i.i.d.}}{\sim} 
\text{Bernoulli}(0.5)$.

\noindent \textbf{Accuracy.} Since each $\mathbf{1}[\cdot]$ is a 
Bernoulli random variable, the standard error of $\hat{\rho}_j^1$ is 
$\sqrt{\rho_j^1(1-\rho_j^1)/M} \leq 1/(2\sqrt{M})$. At 
$M = 10^7$, this upper bound is $0.016\%$---sufficient for 
initialization purposes.

\noindent \textbf{GPU acceleration.} The $M \times n$ matrix of random 
binary samples can be generated and evaluated entirely on a GPU using 
PyTorch tensor operations. On a T4 GPU (Google Colab), computing all 
$n$ probability estimates for $n = 100$, $m = 3$, $M = 10^7$ takes 
5--8 seconds. This is a one-time precomputation before the GA begins.

\noindent \textbf{Failure mode.} When $\hat{\rho}_j^1 = 0$ 
(no Monte Carlo sample is feasible with item $j$ included), 
$p_j = 0$ and item $j$ is never selected. This is the correct behavior 
when item $j$ truly blocks all feasible completions, but can be an 
estimation artifact on very tight instances. In our experiments at 
$\alpha = 0.25$, some items have estimated $p_j = 0$ (never include) 
and others $p_j = 1$ (always include), reflecting genuine extreme 
constraint structure rather than sampling failure.

\subsection{Complexity Analysis}
\label{sec:complexity}

\noindent \textbf{Exact computation.} Expanding each sub-generating 
function $g_j$ enumerates $2^{n-1}$ terms, and checking feasibility 
against $m$ constraints requires $O(m)$ work per term. Since there are 
$n$ items, Phase 1 costs $O(n \cdot 2^n \cdot m)$---exponential in $n$ 
and intractable for $n > 20$ in practice.

\noindent \textbf{Monte Carlo approximation.} Each of the $n$ items 
requires $M$ sample evaluations, each checking $m$ constraints. Phase 1 
therefore costs $O(n \cdot M \cdot m)$---linear in $n$, $M$, and $m$, 
and fully parallelizable across the $M$ sample dimension on a GPU. 
Concretely, for $n = 100$, $m = 3$, $M = 10^7$, this involves 
approximately $3 \times 10^9$ scalar operations, parallelized to 5--8 
seconds on a T4 GPU. The approximation thus reduces an exponential 
enumeration to a polynomial estimation at the cost of a controllable 
error bounded by $1/(2\sqrt{M})$.

Phase 2 (population initialization) costs $O(P \cdot n)$ in both cases 
and is dominated by Phase 1 for any reasonable population size $P$.

\noindent \textbf{Comparison to LP-based approaches.} LP relaxation-based 
initialization methods require solving an optimization subproblem and 
depend on an objective function. The GF Monte Carlo approach requires 
only weights and capacities, solves no optimization problem, and scales 
linearly in $n$, $m$, and $M$.

\noindent \textbf{Memory.} The Monte Carlo approach stores an 
$M \times (n-1)$ binary matrix in GPU memory. At $M = 10^7$ and 
$n = 100$, this is approximately $125$ MB---manageable on modern 
hardware and straightforwardly batchable if memory is limited.

% ============================================================
\section{Worked Examples}
\label{sec:examples}
% ============================================================

\subsection{Single Constraint Example}

Consider the following MKP instance with $n = 3$ items 
and $m = 1$ constraint:

\begin{table}[H]
\centering
\begin{tabular}{lccc}
\toprule
 & Item 1 & Item 2 & Item 3 \\
\midrule
Weight & 3 & 2 & 4 \\
Value  & 5 & 4 & 3 \\
\bottomrule
\end{tabular}
\caption{3-item single constraint instance. Capacity $W = 5$.}
\end{table}

\noindent \textbf{Step 1 --- Full Generating Function}

\begin{flushleft}
\hspace{2em}$f(x) = (1 + x^3)(1 + x^2)(1 + x^4)$
\end{flushleft}
\begin{flushleft}
\hspace{2em}$= 1 + x^2 + x^3 + x^4 + x^5 + x^6 + x^7 + x^9$
\end{flushleft}

Of the $2^3 = 8$ subsets, the feasible ones (exponent $\leq 5$) are: 
$\emptyset$, $\{2\}$, $\{1\}$, $\{3\}$, $\{1,2\}$---giving 
overall feasibility density $5/8 = 0.625$.

\noindent \textbf{Step 2 --- Sub-Generating Functions}

Each $g_j(x)$ is the product over the two items other than $j$ and is 
evaluated against two thresholds: $W - w_j$ (item included) and $W$ 
(item excluded).

\vspace{0.5em}
\noindent \textit{Item 1} $(w_1 = 3,\ W - w_1 = 2)$:
$g_1(x) = (1 + x^2)(1 + x^4) = 1 + x^2 + x^4 + x^6$. 
Of 4 terms: 2 have exponent $\leq 2$ and 3 have exponent $\leq 5$.

\begin{flushleft}
\hspace{2em}$\rho_{1}^{\,1} = 0.50 \quad 
\rho_{1}^{\,0} = 0.75 \quad
p_1 = 0.40$
\end{flushleft}

\noindent \textit{Item 2} $(w_2 = 2,\ W - w_2 = 3)$:
$g_2(x) = (1 + x^3)(1 + x^4) = 1 + x^3 + x^4 + x^7$. 
Of 4 terms: 2 have exponent $\leq 3$ and 3 have exponent $\leq 5$.

\begin{flushleft}
\hspace{2em}$\rho_{2}^{\,1} = 0.50 \quad 
\rho_{2}^{\,0} = 0.75 \quad
p_2 = 0.40$
\end{flushleft}

\noindent \textit{Item 3} $(w_3 = 4,\ W - w_3 = 1)$:
$g_3(x) = (1 + x^3)(1 + x^2) = 1 + x^2 + x^3 + x^5$. 
Of 4 terms: 1 has exponent $\leq 1$ and 4 have exponent $\leq 5$.

\begin{flushleft}
\hspace{2em}$\rho_{3}^{\,1} = 0.25 \quad 
\rho_{3}^{\,0} = 1.00 \quad
p_3 = 0.20$
\end{flushleft}

\noindent \textbf{Step 3 --- Summary and Expected Feasibility Rate}

\begin{table}[H]
\centering
\begin{tabular}{lccc}
\toprule
 & Item 1 & Item 2 & Item 3 \\
\midrule
Weight                          & 3    & 2    & 4    \\
$\rho_{j}^{\,1}$                & 0.50 & 0.50 & 0.25 \\
$\rho_{j}^{\,0}$                & 0.75 & 0.75 & 1.00 \\
\textbf{$p_j$ (ours)}           & \textbf{0.40} & \textbf{0.40} & \textbf{0.20} \\
$p_j$ (uniform)                 & 0.50 & 0.50 & 0.50 \\
\bottomrule
\end{tabular}
\caption{Per-item feasibility densities and sampling probabilities 
for the single-constraint example.}
\end{table}

Item 3 consumes 80\% of the capacity upon inclusion; the framework 
pushes $p_3$ to $0.20$. Enumerating all 8 strings explicitly, the 
expected feasibility rate under our method is $\mathbf{87.2\%}$ 
versus $62.5\%$ under uniform initialization---a gain of $24.7$ 
percentage points.

\subsection{Multiple Constraint Example}

Consider the following MKP instance with $n = 4$ items 
and $m = 2$ constraints:

\begin{table}[H]
\centering
\begin{tabular}{lcccc}
\toprule
 & Item 1 & Item 2 & Item 3 & Item 4 \\
\midrule
Weight (C1) & 2 & 4 & 3 & 5 \\
Weight (C2) & 3 & 2 & 5 & 1 \\
Value       & 4 & 7 & 5 & 3 \\
\bottomrule
\end{tabular}
\caption{4-item two-constraint instance. $W_1 = 7$, $W_2 = 8$.}
\end{table}

\noindent Of the $2^4 = 16$ possible solutions, 9 are feasible (overall 
density $56.25\%$). The two-constraint setting requires evaluating each 
completion against both thresholds simultaneously. We work through 
Item~4 in full to illustrate how the multi-dimensional threshold check 
works, then summarize all four items.

\noindent \textbf{Sub-generating function for Item 4}
$(w_{41} = 5,\ w_{42} = 1)$:

\begin{flushleft}
\hspace{1em}$g_4(x_1, x_2) = (1 + x_1^2 x_2^3)(1 + x_1^4 x_2^2)(1 + x_1^3 x_2^5)$
\end{flushleft}

\noindent When Item 4 is \textbf{included}, the remaining capacity is 
$W_1 - w_{41} = 2$ and $W_2 - w_{42} = 7$. When Item 4 is 
\textbf{excluded}, the full capacity $W_1 = 7$, $W_2 = 8$ is available. 
Table~\ref{tab:g4} enumerates all $2^3 = 8$ completions of $g_4$.

\begin{table}[H]
\centering
\small
\begin{tabular}{lcccc}
\toprule
$\{1,2,3\}$ subset & $e_1$ & $e_2$ & Incl & Excl \\
\midrule
$\emptyset$   & 0 & 0  & \cmark & \cmark \\
$\{3\}$       & 3 & 5  & \xmark & \cmark \\
$\{2\}$       & 4 & 2  & \xmark & \cmark \\
$\{2,3\}$     & 7 & 7  & \xmark & \cmark \\
$\{1\}$       & 2 & 3  & \cmark & \cmark \\
$\{1,3\}$     & 5 & 8  & \xmark & \cmark \\
$\{1,2\}$     & 6 & 5  & \xmark & \cmark \\
$\{1,2,3\}$   & 9 & 10 & \xmark & \xmark \\
\bottomrule
\end{tabular}
\caption{All 8 completions of $g_4$. $e_1$, $e_2$ are total weights 
from items 1--3. Incl: $e_1 \leq 2$, $e_2 \leq 7$. 
Excl: $e_1 \leq 7$, $e_2 \leq 8$.}
\label{tab:g4}
\end{table}

\noindent From Table~\ref{tab:g4}: 2 of 8 completions are feasible 
when Item 4 is included; 7 of 8 are feasible when it is excluded.

\begin{flushleft}
\hspace{1em}$\rho_{4}^{\,1} = 0.250 \quad
\rho_{4}^{\,0} = 0.875 \quad
p_4 = 0.222$
\end{flushleft}

\noindent Item 4 consumes $71\%$ of $W_1$ on its own, leaving almost no 
room for any other item on C1 when it is included---which is why $p_4$ 
is pushed down to $0.222$. Notice that this signal comes specifically 
from C1: Item 4's weight on C2 is only $w_{42} = 1$, which is harmless. 
The two-dimensional threshold check automatically identifies which 
constraint is doing the work.

The same procedure applied to all four items yields:

\begin{table}[H]
\centering
\begin{tabular}{lcccc}
\toprule
 & Item 1 & Item 2 & Item 3 & Item 4 \\
\midrule
Weight (C1)           & 2     & 4     & 3     & 5     \\
Weight (C2)           & 3     & 2     & 5     & 1     \\
$\rho_{j}^{\,1}$      & 0.500 & 0.375 & 0.375 & 0.250 \\
$\rho_{j}^{\,0}$      & 0.625 & 0.750 & 0.750 & 0.875 \\
\textbf{$p_j$ (ours)} & \textbf{0.444} & \textbf{0.333} & \textbf{0.333} & \textbf{0.222} \\
$p_j$ (uniform)       & 0.500 & 0.500 & 0.500 & 0.500 \\
\bottomrule
\end{tabular}
\caption{Per-item probabilities for the two-constraint example.}
\label{tab:multiconstrain}
\end{table}

\noindent Enumerating all 16 strings with these $p_j$ values, the 
expected feasibility rate under our method is $\mathbf{83.8\%}$ versus 
$56.25\%$ under uniform initialization---a gain of $27.6$ percentage 
points.

Note that Items 2 and 3 are penalized equally ($p_j = 0.333$) despite 
very different weight profiles: Item 2 is heavy on C1 and light on C2, 
while Item 3 is the reverse. Both are constrained equally by the 
combination of the two thresholds. This joint accounting---automatically 
identifying which constraints bind for each item---is what distinguishes 
the multi-constraint GF approach from heuristics based on single 
constraints considered independently.

% ============================================================
\section{Experiments}
\label{sec:experiments}
% ============================================================

\subsection{Setup}

We generate a synthetic MKP instance with $n = 100$ items and $m = 3$ 
constraints using a fixed random seed (42). Item weights $w_{ji}$ are 
drawn uniformly from $[1, 100]$ and values $v_j$ uniformly from $[1, 100]$. 
Capacities are set to $W_i = \alpha \cdot \sum_j w_{ji}$ for 
$\alpha \in \{0.25, 0.50, 0.75\}$, giving three instances ranging from 
tight to loose.

We compare two initialization strategies within an otherwise identical 
steady-state GA:

\begin{enumerate}
    \item \textbf{Uniform (baseline):} Each bit set to 1 with probability 
    0.5 independently.
    \item \textbf{GF-Biased:} Each bit $j$ set to 1 with probability 
    $\hat{p}_j$ estimated via Monte Carlo (Section~\ref{sec:montecarlo}).
\end{enumerate}

We use a \textbf{death penalty} fitness model: infeasible solutions 
receive fitness 0 and are non-competing in selection. This isolates the 
effect of initialization cleanly, since no repair mechanism can 
compensate for an all-infeasible initial population.

For the extended experiments at $\alpha = 0.25$ 
(Sections~\ref{sec:hybrid}--\ref{sec:pipeline}), we additionally evaluate 
a hybrid initialization sweep---varying the fraction of GF-biased 
solutions in the initial population---followed by an ILS post-processing 
extension to assess the quality ceiling reachable with downstream 
optimization:

\begin{enumerate}
  \setcounter{enumi}{2}
    \item \textbf{Hybrid:} A mixed initial population combining GF-biased 
    and uniformly initialized solutions at varying proportions (0\%, 25\%, 
    50\%, 75\%, 100\% GF-biased), to identify the minimum informed 
    fraction needed for reliable convergence.
    \item \textbf{GF-Biased + ILS:} GF-biased initialization followed 
    by iterated local search post-processing ($r = 50$ restarts, 
    $k = 5$ bit-flip perturbations), to show the quality ceiling 
    achievable with standard local search on top of the GF output.
\end{enumerate}

\subsection{Parameters}

\begin{table}[H]
\centering
\begin{tabular}{ll}
\toprule
Parameter & Value \\
\midrule
Population size $P$ & 50 \\
Offspring per run & 20,000 \\
Mutation rate & $1/n = 0.01$ \\
Crossover & Uniform \\
Selection & Binary tournament ($k=2$) \\
Fitness & Death penalty \\
Independent runs & 20 \\
Monte Carlo samples $M$ & $10^7$ \\
ILS restarts $r$ & 50 \\
ILS perturbation size $k$ & 5 (bit-flips) \\
\bottomrule
\end{tabular}
\caption{Experimental parameters. ILS parameters apply only to 
Section~\ref{sec:pipeline}.}
\end{table}

\subsection{Reference Best}

A reference best fitness is established per $\alpha$ value by running 
a repair-based GA (Chu \& Beasley) for 30,000 offspring over 20 runs 
and taking the best observed value. These values ($2375$ at $\alpha = 0.25$; 
$3909$ at $\alpha = 0.50$; $4823$ at $\alpha = 0.75$) are not proven 
optimal but serve as a consistent performance ceiling for relative 
comparisons. All quality figures reported below are expressed as a 
percentage of this ceiling. We note that with ILS post-processing, the 
GF-biased approach reaches within $0.5$ percentage points of this 
ceiling on the tight instance (Section~\ref{sec:pipeline}).

% ============================================================
\section{Results}
\label{sec:results}
% ============================================================

\subsection{Sampling Probability Behavior}

Table~\ref{tab:pj} summarizes the distribution of estimated $\hat{p}_j$ 
values across all items at each tightness level. 

\begin{table}[H]
\centering
\begin{tabular}{lcccc}
\toprule
$\alpha$ & Min & Max & Mean & Interpretation \\
\midrule
0.25 & 0.000 & 1.000 & 0.240 & Strong diff. (tight) \\
0.50 & 0.414 & 0.493 & 0.449 & Weak diff. (medium) \\
0.75 & 0.500 & 0.500 & 0.500 & None (loose) \\
\bottomrule
\end{tabular}
\caption{Distribution of $\hat{p}_j$ across all 100 items at each 
tightness level.}
\label{tab:pj}
\end{table}

At $\alpha = 0.75$, every item has $\hat{p}_j = 0.5$ by construction: 
the instance is loose enough that including or excluding any single item 
has negligible effect on feasibility, so the framework degenerates 
automatically to uniform initialization. This is a theoretical prediction 
of the framework (Section~\ref{sec:samplingprob}), and Table~\ref{tab:pj} 
confirms it holds exactly. The identical results at $\alpha = 0.75$ 
reported below are therefore not a comparison between two methods---they 
are a validation that the self-calibrating property holds empirically. 
At $\alpha = 0.25$, the range $[0, 1]$ indicates that some items are 
so weight-dominant that they are never compatible with any feasible 
completion ($p_j = 0$), while others are always compatible ($p_j = 1$).

\subsection{Initialization Feasibility and Final Fitness}

Table~\ref{tab:main} reports initialization feasibility and final 
solution quality jointly. We present them together because at 
$\alpha = 0.25$ the two are causally linked in the most direct sense 
possible: zero initial feasibility under uniform initialization means 
selection never operates, so final fitness is also zero. The table 
tells a single story about what the initialization method enables.

\begin{table}[H]
\centering
\begin{tabular}{llcc}
\toprule
$\alpha$ & Method & Init Feas. & Final (\% ref) \\
\midrule
0.25 & Uniform  & 0.0\%  & 0.0\%  \\
0.25 & GF-Biased & \textbf{52.5\%} & \textbf{96.5\%} \\
\midrule
0.50 & Uniform  & 32.6\% & 98.5\% \\
0.50 & GF-Biased & \textbf{71.6\%} & \textbf{98.8\%} \\
\midrule
0.75 & Uniform  & 100.0\% & 99.5\% \\
0.75 & GF-Biased & 100.0\% & 99.5\% \\
\bottomrule
\end{tabular}
\caption{Initialization feasibility and final fitness under a death 
penalty model. Averages over 20 runs. The $\alpha = 0.75$ rows are 
identical because $p_j = 0.5$ for all items at this tightness level.}
\label{tab:main}
\end{table}

The result at $\alpha = 0.25$ is the central finding. Under uniform 
initialization, every solution in every initial population across all 
20 runs is infeasible. The death penalty assigns fitness 0 to all 
solutions, selection cannot differentiate between them, and no 
improvement ever occurs. GF-biased initialization starts with 52.5\% 
of the initial population feasible, selection operates immediately, 
and the GA converges to 96.5\% of the reference best. This is not a 
marginal improvement---it is the difference between a completely 
inoperative algorithm and one that solves the problem effectively.

At $\alpha = 0.50$, both methods eventually converge to comparable 
final quality (98.5\% versus 98.8\%), but GF-biased initialization 
starts with over twice the fraction of feasible solutions (71.6\% 
versus 32.6\%) and converges faster (Table~\ref{tab:convergence}).

\subsection{Hybrid Initialization}
\label{sec:hybrid}

Before adding post-processing, we first examine how much of the 
population needs to be GF-biased for the algorithm to function at all. 
Table~\ref{tab:hybrid} sweeps the informed fraction at $\alpha = 0.25$.

\begin{table}[H]
\centering
\begin{tabular}{cccc}
\toprule
Inf. Frac. & Init Feas. & Quality (\% ref) & Solved \\
\midrule
0\%   & 0.0\%  & 0.0\%  & 0/20 \\
25\%  & 11.6\% & 97.1\% & 20/20 \\
50\%  & 25.1\% & 96.7\% & 20/20 \\
75\%  & 37.9\% & 97.1\% & 20/20 \\
100\% & 50.5\% & 96.6\% & 18/20 \\
\bottomrule
\end{tabular}
\caption{Hybrid initialization sweep at $\alpha = 0.25$. ``Inf.\ Frac.'' 
is the GF-biased proportion of the initial population; the remainder 
are uniform. ``Solved'' counts runs reaching $\geq 95\%$ of the 
reference best. 20 runs per condition.}
\label{tab:hybrid}
\end{table}

Seeding just 25\% of the population with GF-biased solutions---while 
initializing the remaining 75\% uniformly---is sufficient to rescue 
the algorithm: all 20 runs converge to 97\%+ quality, compared to 0/20 
under pure uniform initialization. Mean solution quality is approximately 
flat across all non-zero informed fractions (96.6--97.1\%), suggesting 
that the critical transition is binary: any non-zero informed fraction 
crosses the feasibility threshold required for selection to operate. 
Notably, the pure GF condition (100\% informed) solves 18/20 runs rather 
than 20/20, while mixed conditions at 25--75\% solve all 20; this is 
consistent with lower population diversity at initialization when all 
solutions are drawn from the same biased distribution. The informed 
fraction is itself a tunable parameter; practitioners can trade off 
initialization overhead against early-generation diversity as their 
application demands.

\subsection{Full Pipeline at $\alpha = 0.25$}
\label{sec:pipeline}

Table~\ref{tab:pipeline} shows what is achievable at $\alpha = 0.25$ 
as downstream optimization is added to the GF-biased baseline.

\begin{table}[H]
\centering
\begin{tabular}{lcc}
\toprule
Method & Quality (\% ref) & Solved \\
\midrule
Uniform + death penalty & 0.0\% & 0/20 \\
GF-Biased & 96.5\% & 20/20 \\
GF-Biased + ILS ($r=50$, $k=5$) & \textbf{99.3\%} & 20/20 \\
\bottomrule
\end{tabular}
\caption{Solution quality at $\alpha = 0.25$ across pipeline 
configurations. 20 runs each. ``Solved'' counts runs reaching 
$\geq 95\%$ of the reference best. A repair-based GA achieves 
$99.8\%$ as a non-comparable reference (it uses objective values 
unavailable to the death-penalty pipeline).}
\label{tab:pipeline}
\end{table}

Adding a lightweight ILS step to the GF-biased GA brings mean quality 
to 99.3\% of the reference best. The ILS configuration used here 
($r = 50$ restarts, $k = 5$ bit-flip perturbations) is one point in a 
tunable parameter space; we do not claim these settings are optimal. 
The point is that the output of the GF framework slots directly into 
standard local search techniques, and substantial quality gains are 
available through straightforward downstream optimization.

\subsection{Convergence Speed}
\label{sec:convergence}

Table~\ref{tab:convergence} reports the median offspring evaluated to 
reach each fitness threshold at $\alpha \in \{0.25, 0.50\}$. The 
$\alpha = 0.75$ case is omitted as both methods are identical by 
construction. Convergence speed is not a primary contribution of this 
work; these results are included for completeness.

\begin{table}[H]
\centering
\small
\begin{tabular}{llccccc}
\toprule
$\alpha$ & Method & 70\% & 80\% & 90\% & 95\% & 98\% \\
\midrule
0.25 & Uniform  & \multicolumn{5}{c}{never (0/20 runs)} \\
0.25 & GF & 481 & 1{,}023 & 2{,}309 & 7{,}529 & 16{,}845$^*$ \\
\midrule
0.50 & Uniform  & 82 & 316 & 893 & 1{,}740 & 8{,}327 \\
0.50 & GF & 54 & 255 & 788 & 1{,}577 & 6{,}724 \\
\bottomrule
\end{tabular}
\caption{Median offspring to first reach each fitness threshold 
(\% of reference best). $^*$reached in only 1 of 20 runs within 
20,000 offspring.}
\label{tab:convergence}
\end{table}

GF-biased initialization also yields modestly faster early convergence 
at $\alpha = 0.50$, consistent with the higher initial feasibility 
giving selection an earlier head start.

% ============================================================
\section{Discussion}
\label{sec:discussion}
% ============================================================

\subsection{When Does GF-Biased Initialization Help?}

The results suggest a clear characterization of when the framework adds 
value. It helps when (1) the instance is tight enough that uniform 
initialization produces non-trivial infeasibility ($\alpha \lesssim 0.50$), 
and (2) the fitness model is constraint-respecting without repair---whether 
that is a death penalty, a soft penalty with a high coefficient, or the 
absence of an objective function entirely. When the instance is loose 
($\alpha = 0.75$), the framework degenerates to uniform initialization 
automatically. This graceful degradation is a theoretical property of 
the $p_j$ formula, not a tuned behavior, and the experiments confirm it.

\subsection{Applicability Beyond Standard MKP}

The framework's most important property is that it is 
\textit{objective-free}: $p_j$ depends only on weights $w_{ji}$ and 
capacities $W_i$, not on values $v_j$. This makes it applicable to 
settings where repair operators are undefined or unreliable:

\begin{itemize}
    \item \textbf{Constraint satisfaction problems (CSPs):} When no 
    objective function exists, the efficiency ranking $u_j$ cannot be 
    computed and repair is undefined. GF-biased initialization provides 
    a principled starting point for any constraint-feasible-only search.
    \item \textbf{Noisy or uncertain objectives:} When value estimates 
    are unreliable, $u_j$ may be misleading. The GF framework ignores 
    objective values entirely and remains well-calibrated.
    \item \textbf{Adversarial correlation structure:} Hill and 
    colleagues~\cite{hill2013neg} showed that standard greedy repair 
    degrades on instances where item weights and values are negatively 
    correlated, because the efficiency ranking systematically 
    disadvantages heavy, high-value items. GF-biased initialization is 
    unaffected by value-weight correlation. Demonstrating this advantage 
    empirically is a primary direction for follow-up work 
    (Section~\ref{sec:futurework}).
\end{itemize}

\subsection{Computational Overhead}

The one-time precomputation at $M = 10^7$, $n = 100$, $m = 3$ runs in 
5--8 seconds on a T4 GPU (Section~\ref{sec:complexity}). For larger-scale 
applications (e.g., $n = 500$ as in OR-Library benchmarks), the cost 
scales as $O(n \cdot M \cdot m)$---linearly in $n$---and is fully 
GPU-parallelizable across the sample dimension.

\noindent \textbf{Amortization across population size.}
The precomputation cost is fixed per instance, independent of population 
size $P$ and offspring budget. The amortized cost per initialized 
solution is therefore $T_{\text{precomp}} / P$: approximately 0.10--0.16 
seconds per solution at $P = 50$, dropping to 0.025--0.04 seconds at 
$P = 200$. Larger populations make the precomputation relatively cheaper 
per solution---the opposite of the repair operator, whose cost is 
incurred once per offspring regardless of population size.

\noindent \textbf{Amortization across runs.}
If the same instance is evaluated $k$ times---for hyperparameter tuning, 
multi-seed experiments, or benchmark evaluation---the $p_j$ vector is 
computed once and reused across all $k$ runs. The amortized precomputation 
cost per run is $T_{\text{precomp}} / k$, which vanishes as $k$ grows. 
This is directly relevant to standard benchmark practice: OR-Library 
evaluations typically run 20--30 trials per instance, reducing the 
effective precomputation overhead to well under one second per trial.

\noindent \textbf{Contrast with repair.}
These cost structures differ fundamentally. The repair operator is applied 
to every infeasible offspring; at $\alpha = 0.25$ under uniform 
initialization, essentially every offspring is infeasible, so repair 
runs on close to 100\% of evaluations throughout the GA budget. Its 
total cost therefore scales linearly with offspring count $T$ and 
problem size $n$. GF precomputation does not scale with $T$ at all: 
regardless of how many offspring are evaluated, the $p_j$ values are 
computed exactly once.

\subsection{Limitations}

\begin{itemize}
    \item \textbf{Order-1 approximation:} The framework estimates 
    per-item feasibility independently using order-1 schemas. For items 
    with strong pairwise interactions, the order-1 approximation may 
    over- or under-estimate the true per-item feasibility contribution. 
    Extending to order-2 schemas is straightforward but increases 
    computation by a factor of $n$.
    \item \textbf{Single-instance generalization:} Our experiments use 
    a single instance per tightness level. Results on OR-Library 
    benchmark instances will be presented in follow-up work.
    \item \textbf{Interaction with repair:} On well-structured standard 
    MKP instances, repair effectively compensates for poor initialization 
    and GF-biased initialization offers little advantage. The benefit is 
    expected to emerge on adversarial or objective-free instances where 
    repair is unavailable.
    \item \textbf{ILS and hybrid parameters:} The ILS configuration 
    ($r = 50$, $k = 5$) and the hybrid informed fraction are unoptimized. 
    The experiments show that these extensions can substantially improve 
    quality; they do not establish optimal settings.
\end{itemize}

% ============================================================
\section{Conclusion and Future Work}
\label{sec:conclusion}
% ============================================================

We have presented a generating function-informed initialization framework 
for genetic algorithms on the Multidimensional Knapsack Problem. The 
framework draws on the same algebraic tradition---encoding combinatorial 
structure as coefficients of formal power series---that Hardy and Ramanujan 
used to analyze integer partitions, and applies it to a new question: 
how much feasible space does each item's inclusion or exclusion preserve? 
By computing per-item feasibility densities from sub-generating functions 
and deriving per-bit sampling probabilities proportional to these densities, 
we obtain an initialization method that is principled, purely 
constraint-driven, and self-calibrating: it degenerates to uniform 
initialization exactly when no constraint structure is present.

The central empirical result is that on tight instances ($\alpha = 0.25$) 
under a death penalty fitness model, GF-biased initialization is the 
only method that functions at all---producing 52.5\% initial feasibility 
and converging to 96.5\% of the reference best, while uniform 
initialization finds no feasible solution in 20 independent runs. 
Adding lightweight ILS post-processing raises this to 99.3\%. Seeding 
as little as 25\% of the population with GF-biased solutions is 
sufficient to rescue the algorithm entirely. On medium instances 
($\alpha = 0.50$), GF-biased initialization converges faster while 
matching final quality. On loose instances ($\alpha = 0.75$), both 
methods are identical by theoretical prediction.

\subsection{Future Work}
\label{sec:futurework}

\begin{enumerate}
    \item \textbf{Adversarial instances (Part 2):} Generating MKP instances 
    following the protocol of Hill et al.~\cite{hill2013neg} with negative 
    value-weight correlation to demonstrate regimes where repair-based 
    approaches degrade while GF-biased initialization remains effective.
    \item \textbf{Constraint satisfaction:} Applying the framework to 
    pure CSP instances where repair operators are undefined.
    \item \textbf{OR-Library benchmarks:} Scaling to standard benchmark 
    instances with $n$ up to 500 and $m$ up to 30. The standard practice 
    of running 20--30 trials per instance makes the amortized precomputation 
    cost per trial negligible.
    \item \textbf{Higher-order schemas:} Extending feasibility density 
    estimation to order-2 schemas to capture pairwise item interactions.
    \item \textbf{Integration with repair:} Investigating whether 
    GF-biased initialization improves the starting point for repair-based 
    GAs on tight benchmark instances.
\end{enumerate}

% ============================================================
\begin{thebibliography}{9}
\bibitem{chu1998genetic}
P.~C. Chu and J.~E. Beasley,
\textit{A Genetic Algorithm for the Multidimensional Knapsack Problem},
Journal of Heuristics, 4(1):63--86, 1998.

\bibitem{beasley1990orlib}
J.~E. Beasley,
\textit{OR-Library: Distributing Test Problems by Electronic Mail},
Journal of the Operational Research Society, 41(11):1069--1072, 1990.

\bibitem{holland1975adaptation}
J.~H. Holland,
\textit{Adaptation in Natural and Artificial Systems},
University of Michigan Press, Ann Arbor, MI, 1975.

\bibitem{wilf1994genfunc}
H.~S. Wilf,
\textit{generatingfunctionology},
2nd edition, Academic Press, 1994.

\bibitem{kellerer2004knapsack}
H.~Kellerer, U.~Pferschy, and D.~Pisinger,
\textit{Knapsack Problems},
Springer, Berlin, 2004.

\bibitem{hill2013neg}
R.~R. Hill, Y.~K. Cho, and J.~T. Moore,
\textit{Problem Generator Algorithms for the Multidimensional Knapsack Problem},
In \textit{Proceedings of the 45th Winter Simulation Conference}, 2013.

\bibitem{hardyramanujan1918}
G.~H. Hardy and S.~Ramanujan,
\textit{Asymptotic Formulae in Combinatory Analysis},
Proceedings of the London Mathematical Society, 2(17):75--115, 1918.
\end{thebibliography}
% ============================================================

\end{document}
